% Article VII — first editorial draft, 10 September 2026.
% Opening integrated with the proofs of revision 06.
% Provenance and conditions: PLAN_RESIDENCIAS_DEMOSTRATIVAS.md.
% Internal references point to proofs included in the body.

\begin{center}\bfseries Abstract\end{center}
\begingroup
\fontsize{10.5}{12}\selectfont
\setlength{\parskip}{2pt plus 1pt minus .5pt}
\noindent
Torsional amplitude is determined as an explicit functional of the
material state through the spinorial action, elimination of contorsion,
and metric variation. A material sector is constructed in which this
amplitude is positive and conserved. Its energetic contribution produces
\(a^{-6}\) dilution and, in the flat homogeneous realization with
\(\Lambda_{\mathrm{eff}}=0\), positive densities, and \(w<1\), a unique
transverse bounce threshold. For dust, an exact solution is obtained for
all time, with finite curvature and Levi--Civita causal geodesics complete
in both directions.

The construction starts from a formal kernel called Triadic Modular
Holography, henceforth HMT. \emph{All Possible Paths}, henceforth APP,
comprises arithmetic paths on \((\mathbb Z/9\mathbb Z)^2\), with
adjacency given by the additive Cayley graph with cardinal generators
and a toroidal cellular realization, distinguishing sum and product
operations and evaluations; TRIT is the unit
of oriented signature \(\{-1,0,+1\}\), determining local regimes; the
\emph{Topological Phase Kernel}, henceforth TPK, composes the dynamics of
transport and memory. The enriched state preserves sheets, orientation,
remainder, carry, and history during phase return. Dimensional Mechanics
realizes its internal publications through physical quantities and
observables, preserving the structural provenance of the couplings.

Gravitational normalization is first developed through the longitudinal
composition \(L=(5759/23040)\alpha^{16}L_*\) and its positive radial
metric. In the realization defined by these operations and by the
identification of the reduced gravitational radius, reciprocity with the
action section fixes \(G=c^3L^2/\hbar_{\rm ret}\) in every admissible
mode. The proof preserves the full archive and establishes the conditions
for static and dynamic reduction of memory. The dimensional choice
\(L_*=1\,\mathrm m\), the longitudinal and metric rules, and the
radial identification are stated in the theorem itself.

Route transport admits a Cartan realization compatible with composition,
inversion, and refinement. Two material actions are distinguished on that
realization: the minimizing extension of the memory screen and the
spinorial realization. Each retains its own variation with respect to
the connection. In the spinorial sector, affine in contorsion, inversion
of the Holst and torsion operators determines the effective correction;
for the minimal extension, the nonlinear dependence of the current is
preserved and a local criterion is proved through the Hessian. The
current, quartic observable, and its normalization explain the value of
the amplitude as well as its role in cosmological evolution.

Persistence of the bounce under perturbations is proved through distinct
local bounds, whereas causal completeness belongs to the explicit global
dust solution. Residual balances, horizon information, and the observer's
operations preserve the conditions on sources, couplings, and charts
that they use. The relational response of the joint state is connected
to advanced and retarded propagation through Green operators for finitely
supported sources, associated with the covariant difference constructed
from route transport and a boundary
compatibility condition. The result relates generated memory, its
geometric realization, and the dynamical and informational consequences
of that realization.

The composition of action, the Barbero--Immirzi area increment, and the
horizon determines the conversion
\(T_{\rm cos}(L_\alpha)/\Theta_{\rm clk}=\log3/(2\pi)\).
The angular response Hamiltonian yields a finite free-energy balance
through relative entropy and a first law at fixed radius. A relative
thermometric section also characterizes the positive transducer and
the criterion for identifying another reader on that same section.
\par\endgroup
\clearpage

\section{Introduction}
\label{vii:sec:introduccion}

The gravitational question organizing this work is how a structure
preserving the history of its operations can determine a material current
and, through it, a contribution to spacetime evolution. The answer is
constructed from route transport through to variation of the action:
the central result is an explicit torsional amplitude, accompanied by a
sector of conserved positivity and a global cosmological solution in the
dust realization. The thesis connects arithmetic memory to a geometric
response determined by operations; phase return preserves an observable
periodicity while the state accumulates orientation, carry, and depth.

A formal mathematical kernel called Triadic Modular Holography,
henceforth HMT, is introduced for this purpose. \emph{All Possible
Paths}, henceforth APP, is an arithmetic path structure on
\((\mathbb Z/9\mathbb Z)^2\). Its adjacency is described by the additive
Cayley graph with cardinal generators \(\pm(1,0)\) and \(\pm(0,1)\),
and its periodic cellular realization is toroidal. Additive
accumulation, multiplicative composition, and their evaluations retain
distinct roles within that structure. TRIT is the unit of oriented
signature \(\{-1,0,+1\}\), typing an operation's local regime. The
\emph{Topological Phase Kernel}, henceforth TPK, constitutes its dynamics
of selection, transport, memory, and prolongation. The enriched state and
the joint discrete structure of the continuum preserve sheets and
residual data through those operations. Dimensional Mechanics realizes
their outputs in action sections, geometric quantities, and physical
observables; internal values transmit to each stage the relations
determining them.

The central question is how transported memory comes to enter
gravitational evolution. Dilution, the bounce, and its persistence
describe the cosmological consequences of an amplitude that must first
be constructed from the material action. The first relationship is expressed
through two readers of one and the same route,
\begin{equation}
 \frac{a_n}{a_{\mathrm{ref}}}=\varphi^{n/3},
 \qquad M_n=\varphi^{-2n},
 \qquad
 M_n=\left(\frac{a_n}{a_{\mathrm{ref}}}\right)^{-6}.
 \label{vii:eq:programa-dilucion}
\end{equation}
Here \(a_{\mathrm{ref}}>0\) fixes the normalization of the scale factor.
The power expresses the relationship between the growth of the spatial
realization and the quadratic character of the memory reader. The gravitational
contribution also requires determining the current and amplitude accompanying
that scale dependence. Both constructions occupy a section of their own
before being introduced into the cosmological equations.

The first movement of the article brings together construction and geometry.
The decomposition of transport distinguishes symmetric memory from oriented
circulation. The representation of the route groupoid in Cartan geometry
is developed with composition, inversion, covariance, and refinement.
The action, the Barbero--Immirzi parameter, the gravitational constant, and the
areal unit preserve their structural antecedents; their reuse
transmits the conditions individuating each section.

In particular, Section~\ref{g26:sec:rigidez-radial-en} first constructs
the length \(L\), composes its positive reader with the same character
that determines energy, and obtains the gravitational coefficient.
The circular clock then receives
\(t_0=\pi_{\rm HMT}L/(54c)\); its radius \(R_{108}\) coincides
with \(L\) in this realization. The proof distinguishes longitudinal
composition, operator response, and evaluation in physical units, in that
order. The longitudinal and metric rules are compositions made explicit
in this development; the theorem proves rigidity for the realizations
preserving them, with the stated radial identification.

The second movement develops amplitude, evolution, and persistence.
The screen current is carried to a differential with respect to the
connection through transported incidence and the minimizing extension.
The spinorial realization is developed, in turn, from the internal
Clifford matrices and its material action. The minimizing-extension
action and the spinorial action are distinct functionals: the former
generally depends nonlinearly on the connection, whereas the latter is
affine in contorsion with cotetrad and spinorial field fixed. Their
variations act on each functional's own arguments; any composition of
the two must preserve that dependence.
Tracial normalization
fixes the corresponding density reading and remains distinguished from
an extensive sum. Variation of the connection determines the torsional
response of the minimal action. As the cotetrad and material fields evolve,
the torsion algebraically associated with their spin current also changes.

The amplitude of the spinorial sector is determined by
\begin{equation}
 s_0[\varrho,V_c]
 =\frac{3\kappa c^2\hbar^2}{16}
   \frac{\gamma^2}{1+\gamma^2}
   \frac{\operatorname{Tr}
    (\varrho\widehat{\mathscr Q}_{\rm sp})}{V_c^2},
 \qquad \kappa=\frac{8\pi G}{c^4}.
 \label{vii:eq:programa-amplitud}
\end{equation}
Here \(\varrho\) is the normalized material state, \(V_c\) the comoving
volume, and \(\widehat{\mathscr Q}_{\rm sp}\) the quartic observable
constructed from the current.
Theorem~\ref{vii:thm:amplitud-material-explicita} derives the formula, and
Corollary~\ref{vii:cor:dominio-positivo-conservado} exhibits states
in which the amplitude is positive and remains constant. Dependence
on the second moment preserves information about the state even when
the mean current vanishes. The material datum and its volume fix a
realization; the coupling constants come from the structural
sections constructed previously.

Elimination of torsion is also treated for the minimal-extension
action, whose dependence on the connection is generally
nonlinear. Its differential determines the stationary equation;
the Hessian controls local continuation, and an integral identity
determines the effective action. This formulation preserves the
full dependence of the minimum when variations are performed.

The bounce formulation preserves the positive densities, the sign of
the torsional contribution, and the exponent inequality of the homogeneous
realization. Its local persistence is proved through two separate
controls: a uniform perturbation bound preserves the sign change
at the endpoints of an interval; a uniform bound on its derivative
preserves monotonicity and transversality. Realizations of the cosmological
selectors are stated with their own domains, together with the
invariants and balances they produce.

The flat dust solution with \(\Lambda_{\mathrm{eff}}=0\)
also provides a global development:
\[
 a(t)=a_{\rm b}\left(1+
       \frac{(t-t_{\rm b})^2}{\tau^2}\right)^{1/3},
 \qquad
 a_{\rm b}^3=\frac{s_0}{\rho_0},\qquad
 \tau^2=\frac{s_0}{6\pi G\rho_0^2},
 \quad c=1.
\]
The proof establishes curvature regularity and completeness of
Levi--Civita causal geodesics in that homogeneous realization.
Its integration and local persistence of the threshold describe two
complementary properties, with their respective explicit domains.

Conservation of the residual tensor is formulated in the selected
metric chart, with constant couplings on that chart and a
covariantly conserved material source. Its decomposition into trace
and traceless parts identifies exchange between the two sectors;
constant trace characterizes the case of separate conservation.

The third movement studies horizons and relational response. The reduced
state, accessible records, and purification describe
what information a reading publishes and what information remains in
the correlations of the joint state. The operational definition of the
observer comprises orientation and trivialization of transport,
restriction to observables, and projection onto the subsystem. The
self-inertial response is connected to those operations through factorization
through the global trace and the geometric projection defect.

\Needspace{24\baselineskip}
Horizon balances likewise preserve the conventions of
area, energy, temperature, and state that they use.
Sections~\ref{viii:subsec:conversion-termica-calendario-horizonte}--\ref{viii:subsec:seccion-termometrica-registro}
compose the area quantum with the calendar's decision energy and obtain
the conversion between chronological and cosmological temperatures.
The reduced state preserves the occupations of both angular sectors.
At fixed radius, its energetic response has Hamiltonian
\(\mathsf H_R=\tau_R\Delta_0\mathsf D_\partial\); its variations
give the first law, while finite differences retain the relative-entropy
term. This development constructs a relative thermometric section and
establishes on it a test of equality between linear conversion maps.
Sectorial provenance remains accessible in both area and energetic cost.
Advanced and retarded propagation for finitely supported sources is
composed on unitary route transport. Equality of their boundary readings is
expressed through a linear constraint on sources; its projection
determines the variational response in that domain.

The argument's unity lies in the relationship between source, current,
amplitude, and geometry. Memory preserves information not exhausted by
the phase reading; its material realization produces a second moment
capable of contributing to the dynamics even with zero mean current.
Cosmological evolution expresses that contribution globally in the
resolved sector, and the observer's operations specify which part of the
information remains accessible. Structure, equation, and solution are
thereby connected within realizations with explicit domains.
